(1) 设实电荷到场点距离 \(r_+=\sqrt{\rho^2+(z-d)^2}\)，镜像点距离 \(r_-=\sqrt{\rho^2+(z+d)^2}\)。上半空间尝试 \(\phi_1=(q/r_++q'/r_-)/(4\pi\varepsilon_1)\)；下半空间尝试 \(\phi_2=q''/(4\pi\varepsilon_2r_+)\)。此处下半空间的 \(r_+\) 表示位于域外实电荷位置的等效源，其在 \(z<0\) 不产生奇点。

(2) 在 \(z=0\)，电势连续给 \((q+q')/\varepsilon_1=q''/\varepsilon_2\)；法向位移连续 \(\varepsilon_1\partial_z\phi_1=\varepsilon_2\partial_z\phi_2\) 给 \(q-q'=q''\)。联立得到 \(q'=q(\varepsilon_1-\varepsilon_2)/(\varepsilon_1+\varepsilon_2)\)，\(q''=2\varepsilon_2q/(\varepsilon_1+\varepsilon_2)\)。解满足远场衰减和各域 Laplace 方程，且边界条件确定唯一静电解。

(3) 求力只计算诱导界面电场，不能包括实电荷自身场。上方镜像距离为 \(2d\)，故 \(F_z=q q'/(4\pi\varepsilon_1(2d)^2)=q^2(\varepsilon_1-\varepsilon_2)/[16\pi\varepsilon_1(\varepsilon_1+\varepsilon_2)d^2]\)。若下方介电常数更大，则力为负（趋向界面）；反之排斥。缓慢将电荷移到距离 \(d\) 所储存的相互作用能 \(U(d)=q q'/(16\pi\varepsilon_1d)\)，满足 \(F_z=-dU/dd\)。

(4) 当 \(\varepsilon_2\to\varepsilon_1\)，\(q'\to0\)，反作用力消失，界面变得不可区分。当 \(\varepsilon_2\to\infty\)，\(q'\to-q\)，力退化为理想接地导体平面 \(F_z=-q^2/(16\pi\varepsilon_1d^2)\)。对于一般有限 \(\varepsilon_2\)，不能把下方介质当作导体，也不能把极化等效成真实自由面电荷；镜像电荷仅是满足边界条件的数学表示。
